% Synthos: A Self-Evolving Cognitive Operating System for AI-Augmented Scientific Research
\documentclass[review]{elsarticle}

\usepackage{amsmath,amssymb}
\usepackage{graphicx}
\usepackage{booktabs}
\usepackage{hyperref}
\usepackage{xcolor}
\usepackage{textcomp}
\usepackage[utf8]{inputenc}
\usepackage{array}

\journal{Intelligent Computing}

\begin{document}

\begin{frontmatter}

\title{Synthos: A Self-Evolving Cognitive Operating System for AI-Augmented Scientific Research}

\author[1]{Xiaokai Yang\textsuperscript{*}}
\ead{yakeworld@outlook.com}
\author[2]{Synthos Project}
\address[1]{Department of Neurology, Wenzhou People's Hospital, Wenzhou 325000, China}
\address[2]{GitHub: \url{https://github.com/yakeworld/Synthos}}

\begin{abstract}
Large language model (LLM) agents have demonstrated remarkable potential in scientific research assistance, yet existing systems remain fragmented between tool orchestration and genuine cognitive reasoning. We present Synthos, a self-evolving cognitive operating system that reifies eight epistemological principles as executable SKILL.md cognitive atoms. Unlike conventional research agents that wrap Python libraries as tools, Synthos operates as a purely skill-driven architecture: the Agent itself serves as the runtime, with zero Python infrastructure code. The system's core innovation lies in its constitutional design---a formal hierarchy of immutable principles---and its self-evolution engine (v2.11), an 11-step state machine with structural probes, functional benchmarks, external absorption scanning, SEPL rollback protocols, and Git-as-Memory versioning. Over 41 evolution cycles, Synthos achieved a composite quality score of 0.975 (EXCELLENT), maintained all cognitive atoms at perfect structural and benchmark scores of 1.00, absorbed patterns from 12 external projects, and produced verified paper drafts across four biomedical domains (biomechanics, ophthalmology, otoneurology, neurology).
\end{abstract}

\begin{keyword}
Cognitive Architecture \sep Self-Evolution \sep Scientific AI Agent \sep Constitutional AI \sep Epistemological Computing
\end{keyword}
\end{frontmatter}

\section{Introduction}
\label{sec:introduction}

\subsection{The Rise of AI Research Agents}
\label{sec:rise}

The pursuit of autonomous scientific discovery through artificial intelligence has produced a diverse ecosystem of research agents~\cite{elovic2024gpt,sakana2025ai,shao2024storm,futurehouse2025paperqa2,asai2025openscholar}. These systems have achieved remarkable benchmarks---PaperQA2 achieving 99.3\% on PubMedQA, AI Scientist producing the first AI-written workshop paper---yet they share a fundamental limitation: their cognitive architecture remains implicit, embedded in orchestration code rather than explicitly modeled as a reasoning framework.

\subsection{Three Critical Gaps}
\label{sec:gaps}

\textbf{Epistemological rigor is weakly enforced.} Most agents lack formal mechanisms for evidence traceability, falsifiability constraints, or systematic bias detection. The SKILL.md paradigm~\cite{hermes2025hermes} introduced executable skill files, but its application to scientific reasoning remains nascent. While cognitive architectures for language agents have been surveyed~\cite{sumers2023cognitive} and generative agent frameworks proposed~\cite{park2023generative}, comprehensive agent surveys~\cite{wang2024surveyagent,xi2023rise} confirm that none formally encode epistemological constraints.

\textbf{Self-improvement remains ad-hoc.} Evolution in existing systems is typically limited to tool-patch updates rather than architectural self-reflection. LangGraph~\cite{langchain2024langgraph} introduced state machines for agent workflows, LangChain~\cite{chase2022langchain} pioneered composable LLM applications, and DSPy~\cite{khattab2023dspy} demonstrated automated prompt optimization. However, none address the challenge of a system improving its own cognitive architecture through a deliberate, self-directed evolution cycle.

\textbf{Philosophical principles remain aspirational.} Even systems that invoke the ``scientific method'' do not execute epistemological constraints as first-class architectural components. Foundational LLM advances---Transformers~\cite{vaswani2017attention}, GPT-3~\cite{brown2020language}, chain-of-thought reasoning~\cite{wei2022chain}, self-consistency~\cite{wang2023selfconsistency}, tree-of-thoughts~\cite{yao2023tree}, and retrieval-augmented generation~\cite{lewis2020retrieval}---have enabled increasingly capable agents, yet the reasoning frameworks that govern these capabilities remain implicit and non-enforceable.

\subsection{Related Work and Positioning}
\label{sec:related}

Recent advances address these gaps from complementary but incomplete directions. Claude Code~\cite{anthropic2025claude} introduced constitutional AI with drift detection, restricting its scope to code generation. DeepResearchAgent~\cite{skyworkai2026deep} proposed SEPL rollback protocols for self-evolution but remains Python-dependent. AI-Research-SKILLs~\cite{orchestra2026airesearch} demonstrated 90+ SKILL.md files with ARA provenance, yet lacks a self-evolution mechanism. GEPA~\cite{gepa2025gepa} introduced reflective strategy selection without constitutional governance. STORM~\cite{shao2024storm} excels at literature synthesis but does not extend to hypothesis generation or verification. OpenScholar~\cite{asai2025openscholar} synthesizes literature at scale but operates as a monolithic pipeline. Nature-skills~\cite{yuan2025nature} and Academic Research Skills~\cite{ars2025academic} advanced the SKILL.md paradigm for task-specific capabilities, while PaperOrchestra~\cite{song2026paperorchestra} demonstrated multi-agent pipeline advantages. AutoR~\cite{autox2025autor} shares the human-in-loop philosophy but implements it via Python harness. KILO-KIT~\cite{vo2025kilo} contributed CRISP-DM templates for structured research workflows. However, no existing system simultaneously achieves explicit epistemological encoding, constitutional governance, self-evolving architecture, and zero-Python execution.

\subsection{Our Contributions}
\label{sec:contributions}

\textbf{Contribution 1: Epistemological Code} --- Eight philosophical frameworks from the history and philosophy of science are reified as executable SKILL.md cognitive atoms with formal I/O contracts, building on representation learning principles~\cite{bengio2013representation}. First Principles governs decomposition; Falsificationism ensures every hypothesis generates explicit counterargument paths; Bayesian Reasoning drives confidence propagation across the verification chain.

\textbf{Contribution 2: Constitutional Hierarchy} --- All system behavior is governed by CONSTITUTION.md v5.0~\cite{yang2026synthosconst}, a formal hierarchy of immutable principles ranked CON~$>$~MEM~$>$~CMD~$>$~SKL~$>$~DEF, enforced by a philosophical immune system.

\textbf{Contribution 3: Self-Evolution Infrastructure} --- An 11-step evolution engine~\cite{yang2026synthos} (v2.11) incorporating PW-Bench~\cite{song2026paperorchestra}, SEPL rollback~\cite{skyworkai2026deep}, Git-as-Memory~\cite{haiyang2025aris}, hypothesis-first single-metric improvement, and active inference gates.

Table~\ref{tab:comparison} compares Synthos with representative existing systems across key dimensions.

\begin{table}[h]\centering\caption{Comparison with representative research agent systems.}\label{tab:comparison}
\small
\begin{tabular}{lcccc}\toprule
\textbf{System} & \textbf{Epistemic} & \textbf{Const.} & \textbf{Self-Evol.} & \textbf{Zero-Py}\\\midrule
GPT-Researcher~\cite{elovic2024gpt} & $\times$ & $\times$ & $\times$ & $\times$\\
AI Scientist~\cite{sakana2025ai} & $\times$ & $\times$ & $\bigtriangleup$ & $\times$\\
PaperQA2~\cite{futurehouse2025paperqa2} & $\times$ & $\times$ & $\times$ & $\times$\\
Claude Code~\cite{anthropic2025claude} & $\times$ & $\checkmark$ & $\times$ & $\times$\\ 
DSPy~\cite{khattab2023dspy} & $\times$ & $\times$ & $\checkmark$ & $\times$\\
DeepResearchAgent~\cite{skyworkai2026deep} & $\times$ & $\times$ & $\checkmark$ & $\times$\\
AI-Research-SKILLs~\cite{orchestra2026airesearch} & $\bigtriangleup$ & $\times$ & $\times$ & $\checkmark$\\
\textbf{Synthos (ours)} & $\checkmark$ & $\checkmark$ & $\checkmark$ & $\checkmark$\\\bottomrule
\end{tabular}
\end{table}

\section{Architecture}

\subsection{Design Principles}

Synthos is built on four immutable principles~\cite{yang2026synthosconst}: (1) Carbon-Silicon Symbiosis---human-in-the-loop; (2) Self-Evolution as First Principle---all outputs serve evolution; (3) Precise Atomic Boundaries---non-overlapping cognitive atoms; (4) Epistemological Honesty Over Sycophancy---claims cite sources. These form a priority hierarchy: CON~$\gg$~MEM~$\gg$~CMD~$\gg$~SKL~$\gg$~DEF, enforced by a philosophical immune system.

\subsection{Cognitive Atoms}

The eight-dimensional cognitive framework maps to six SKILL.md atoms plus a task router, as illustrated in Figure~\ref{fig:architecture}:

\begin{figure}[h]\centering
\includegraphics[width=\textwidth]{fig_architecture.pdf}
\caption{Synthos architecture. Top: constitutional governance layer (CONSTITUTION.md v5.0). Middle: six cognitive atoms (ACQ$\rightarrow$EXT$\rightarrow$ASC$\rightarrow$HYP$\rightarrow$ARG/VER) forming a directed graph with viewpoint-verification feedback loop. The Task Router selects the shortest execution path based on query complexity. Bottom: the 11-step evolution engine operates as a persistent state machine with five conditional gates and three execution paths.}
\label{fig:architecture}
\end{figure}

\begin{table}[h]\centering\caption{Cognitive atoms with epistemological bases and Cycle~41 scores.}\label{tab:atoms}
\small
\begin{tabular}{l>{\raggedright\arraybackslash}p{2.5cm}>{\raggedright\arraybackslash}p{4.5cm}c}\toprule
\textbf{Atom} & \textbf{Epistemic Basis} & \textbf{Core Function} & \textbf{Score}\\\midrule
Task Router & Occam's Razor & Shortest-path routing & 1.00\\
Knowledge Acq. & First Principles & Multi-source paper retrieval & 1.00\\
Knowledge Extr. & First+Occam & Structured extraction & 1.00\\
Association Disc. & Systems Thinking & Cross-paper relation detection & 1.00\\
Hypothesis Gen. & Falsification+Analogy & Testable hypotheses with counterexample paths & 1.00\\
Argument Expr. & Analogy+Sys.Think. & IMRaD academic writing & 1.00\\
Viewpoint Verif. & Bayesian+Falsification & Multi-angle counterargument gen. & 1.00\\\bottomrule
\end{tabular}\end{table}

Each atom is a declarative SKILL.md file with YAML frontmatter and markdown body. No Python infrastructure or orchestration scripts exist---the Agent is the runtime.

\subsection{The Viewpoint Verification Atom: A Case Study in Epistemic Depth}

The \texttt{viewpoint-verification} atom (v1.3.0) exemplifies Synthos's multi-layered epistemic architecture. Its confidence calibration pipeline incorporates three independent layers:

\textbf{Layer~1 (Bayesian)}: Prior confidence starts at 0.70 and is adjusted by counterargument strength (linear deduction 0.05--0.15 per argument), internal weakness penalties (0.08 each), and robustness concern penalties (0.05 each).

\textbf{Layer~2 (Citation Quality)}: Inspired by PaperOrchestra's Citation F1 framework~\cite{song2026paperorchestra}, each reference is classified as P0 (essential) or P1 (supporting), and confidence is modulated by citation quality factor.

\textbf{Layer~3 (Epistemic Scoring)}: Drawing from ARA Rigor Reviewer~\cite{orchestra2026airesearch}, six dimensions---Evidence Relevance, Falsifiability Quality, Scope Calibration, Argument Coherence, Exploration Integrity, and Methodological Rigor---are each scored 1--5 and averaged into an epistemic factor that acts as a final confidence gate.

This three-layer cascade transforms verification from a heuristic judgment into a formally auditable process.

\subsection{Evolution Engine v2.11}

The engine operates as an 11-step state machine:

\begin{verbatim}
Path A: LOAD_CONST -> DRIFT_CHECK -> EXIT
Path B: LOAD_CONST -> LOAD_STATE -> LESSONS ->
  DRIFT_CHECK -> PROBE -> BENCHMARK ->
  [fail?] OPTIMIZE -> [>6h?] EXTERNAL ->
  DIAGNOSE -> [discovery?] INTERRUPT ->
  IMPROVE -> VERIFY -> RECORD
Path C: LOAD_CONST -> restore checkpoint -> continue
\end{verbatim}

Key innovations include \textbf{Hypothesis-First Improvement} (each IMPROVE begins with a falsifiable hypothesis written to \texttt{hypothesis\_\{cycle\}.yaml}, inspired by~\cite{haiyang2025aris}), \textbf{Git-as-Memory} (each cycle creates a branch; passes are kept, failures reverted; agent reads \texttt{git log} before DIAGNOSE), \textbf{Single-Metric Focus} (each cycle improves exactly one dimension; three failed cycles trigger stuck protocol), and \textbf{SEPL Rollback}~\cite{skyworkai2026deep} (OPTIMIZE patches undergo commit gate; VERIFY failure triggers automatic \texttt{git revert}; three consecutive rollbacks lock the skill).

\textbf{Active Inference Gate}: When DIAGNOSE detects high epistemic uncertainty---confidence variance $>0.3$, benchmark variance $>0.2$---the engine generates formal experiment proposals to minimize expected surprise, operationalizing the Free Energy Principle~\cite{friston2010free} at the system architecture level.

\textbf{ARA-Structured Provenance}: Evolution logs classify events into seven types---Decision, Experiment, Dead End, Pivot, Claim, Heuristic, Observation---each with provenance markers, following the ARA framework~\cite{orchestra2026airesearch}.

\begin{figure}[h]\centering
\includegraphics[width=0.85\textwidth]{fig_gates.pdf}
\caption{Quality gate architecture for the paper-writing pipeline. G1--G7 form a sequential verification chain; each gate must pass before the next can proceed. G7 triggers the SCI 7-dimension content review.}
\label{fig:gates}
\end{figure}

\subsection{External Absorption Pipeline}

Absorption scans GitHub ($\ge50$ stars, updated within 1~year), Hermes skills, and academic literature. Projects are evaluated on five dimensions (complementarity, code quality, community activity, integration cost, license compatibility). Scores $\ge4.0/5$ proceed through P0 reconnaissance~$\rightarrow$~P1 core injection~$\rightarrow$~P2 schema integration~$\rightarrow$~P3 meta-cognitive enhancement.

\section{Results}

\subsection{Evolution Performance}

Over 41 cycles (May~11--19,~2026), Synthos demonstrated consistent improvement:

\begin{table}[h]\centering\caption{Evolution metrics across representative cycles.}\label{tab:evolution}
\small
\begin{tabular}{lcccc}\toprule
\textbf{Metric} & \textbf{Cycle~1} & \textbf{Cycle~10} & \textbf{Cycle~23} & \textbf{Cycle~41}\\\midrule
Structural Average & 0.861 & 1.00 & 1.00 & 1.00\\
Benchmark Score & 0.66 & 1.00 & 1.00 & 1.00\\
Composite Score & 0.86 & 0.94 & 0.97 & 0.975\\
Degraded Atoms & 0 & 0 & 0 & 0\\
API Health & degraded (S2~429) & healthy & healthy & healthy\\
External Absorptions & 0 & 1 & 3 & 10\\
Constitution Version & v2.0 & v3.0 & v5.0 & v5.0\\
Evolution Log Lines & 12 & 240 & 480 & 1,313\\\bottomrule
\end{tabular}\end{table}

The composite score stabilized at 0.975 (EXCELLENT) by Cycle~41. No atom has ever been flagged as degraded. Figure~\ref{fig:evolution} shows the full evolution trajectory across 41 cycles.

\begin{figure}[h]\centering
\includegraphics[width=0.85\textwidth]{fig_evolution.pdf}
\caption{Evolution trajectory across 41 cycles. The composite score rose from 0.86 (Cycle~1) to 0.975 (Cycle~41), with structural and benchmark scores achieving 1.00 from Cycle~2 onward.}
\label{fig:evolution}
\end{figure}

\subsection{Benchmark Performance}

Across 7 benchmark suites (ROUTE-01/02/03, ACQ-01/02, EXT-01/02, ASC-01/02, HYP-01/02, ARG-01/02, VER-01/02) and 7 golden I/O contract tests, Cycle~41 achieved 100\% pass rate (14/14). The golden test suite validated schema completeness and anti-sycophancy protocol integrity.

\subsection{External Absorption Record}

\begin{table}[h]\centering\caption{External absorption history.}\label{tab:absorption}
\footnotesize
\begin{tabular}{l>{\raggedright\arraybackslash}p{3.2cm}>{\raggedright\arraybackslash}p{3.5cm}c}\toprule
\textbf{Source} & \textbf{Absorbed Pattern} & \textbf{Target} & \textbf{Score}\\\midrule
Claude Code~\cite{anthropic2025claude} & Hook events, drift detection & Evolution v2.4--v2.5 & --\\
DSPy~\cite{khattab2023dspy} & Type signatures, auto-optimizer & All atoms, Evolution v2.6 & --\\
LangGraph~\cite{langchain2024langgraph} & State machine, checkpoints & Evolution v2.7 & --\\
GEPA~\cite{gepa2025gepa} & Reflective strategy selection & Evolution v2.8 & --\\
OpenAI Agents SDK~\cite{openai2025agents} & Input guardrails & Evolution v2.9 & --\\
AutoResearchClaw~\cite{aiming2025claw} & Citation verification, LaTeX & Extended skills & 4.5\\
KILO-KIT~\cite{vo2025kilo} & CRISP-DM templates & Extended skills & 5.0\\
ARS~\cite{ars2025academic} & Anti-sycophancy protocol, data levels & VER v1.2, CONST. P6 & 4.5\\
PaperOrchestra~\cite{song2026paperorchestra} & Citation F1, LitReview 6-axis & VER v1.2, BENCHMARKS & 4.0\\
DeepResearchAgent~\cite{skyworkai2026deep} & SEPL rollback, auditable lineage & Evolution v2.10 & 4.2\\
AI-Res. SKILLs~\cite{orchestra2026airesearch} & ARA provenance, epistemic scoring & VER v1.3, Evolution v2.10 & 4.5\\
ARIS~\cite{haiyang2025aris} & Git-as-Memory, hypothesis-first & Evolution v2.11 & 4.0\\\bottomrule
\end{tabular}\end{table}

\subsection{Practical Validation}

Synthos has produced paper drafts across four domains: VOR-based 3D Kappa angle estimation (biomechanics, 14~pages, 30~references), iris anatomical segmentation (ophthalmology, 18~pages), BPPV rehabilitation simulation (otoneurology, 12~pages), eye-tracking ADHD screening (neurology, 70+ literature synthesis), and 4 NSFC grant proposals (vertigo biomechanics, fungal sinusitis, PD aspiration prediction, ADHD eye-tracking screening). Each paper follows the pipeline: NotebookLM query extraction~$\rightarrow$~IMRaD LaTeX construction~$\rightarrow$~7-pathway PDF retrieval~$\rightarrow$~quality gate G1-G7~$\rightarrow$~SCI 7-dimension review. The pipeline requires 2--3 human-in-the-loop sessions per paper. Published outputs include a completed manuscript on CRISP-DM Helix data isolation (submitted to Kaggle, featured on the platform homepage) and a Synthos-inspired cognitive framework paper targeting BMC Medical Education.

\section{Discussion}

\subsection{Key Insight: Epistemology as Computational Constraint}

The key insight validated by Synthos is that epistemology is not an abstract philosophical overlay but a computational constraint that can be operationalized at the architecture level. The Free Energy Principle~\cite{friston2010free} drives active inference gates; Model-Dependent Realism manifests as persona-based debate mechanisms within the Viewpoint Verification atom; Falsificationism is operationalized as mandatory counterargument generation with explicit counterexample paths in every Hypothesis Generation cycle. This differs fundamentally from conventional agents that hardcode heuristic quality checks rather than encoding epistemic principles as first-class architectural constraints.

\subsection{Evolution as a Viable Architectural Principle}

The evolution engine's trajectory---structural score from 0.861 to 1.00, composite score from 0.86 to 0.975 over 41 cycles, zero degraded atoms---demonstrates that self-evolution is not merely a theoretical aspiration but a practically viable architectural principle for AI research agents. Each of the 12 external absorptions (Table~\ref{tab:absorption}) improved a specific dimension without degrading existing capabilities, validating the SEPL rollback and Git-as-Memory mechanisms. Recent surveys confirm that LLM-driven scientific discovery~\cite{gao2024scientific} and agent architectures~\cite{wang2024surveyagent,xi2023rise} remain open challenges that Synthos directly addresses.

\subsection{Zero-Python Architecture as a Paradigm Shift}

The finding that a complete research workflow---7 cognitive atoms, router, evolution engine, absorption pipeline---can execute with zero Python infrastructure code has implications beyond Synthos itself. The SKILL.md paradigm, validated across nature-skills~(5,625~stars), AI-Research-SKILLs~(6,358~stars), and PaperOrchestra~(461~stars), demonstrates that declarative skill files are a viable alternative to Python orchestration for agent capabilities. This reduces the barrier to contribution from software engineering to domain expertise: a neurologist extending Synthos writes Markdown, not Python code.

\subsection{Limitations and Future Work}
\label{sec:limitations}

Seven limitations merit acknowledgment. \textbf{(1) Absorption requires human oversight.} While the external absorption pipeline autonomously identifies and evaluates candidate projects, the final integration---SKILL.md adaptation, schema alignment, and I/O contract verification---requires manual intervention. Fully autonomous absorption, where the system rewrites its own SKILL.md files without human review, remains a target for future evolution cycles.

\textbf{(2) Active inference generates proposals without execution.} The Free Energy Principle gate (Section~3.2) successfully identifies high-uncertainty states and generates experiment proposals, but these proposals are recorded rather than autonomously executed. Closing this loop---auto-generated experiments running within the evolution cycle---would substantially accelerate epistemic convergence.

\textbf{(3) Single-runtime dependency.} Synthos currently executes exclusively on Hermes Agent. While the SKILL.md format is theoretically platform-agnostic, cross-platform validation (Claude Code, Cursor, OpenCode) has not been performed. Platform-specific features (file I/O conventions, tool availability, context window management) may require adaptations.

\textbf{(4) Reasoning quality depends on the underlying LLM.} Epistemic rigor is fundamentally bounded by the model's reasoning capabilities. Multi-model epistemic triangulation---running the same cognitive atom across different LLMs and comparing outputs---could mitigate model-specific biases but adds latency and cost.

\textbf{(5) No multi-agent collaboration.} The current architecture assumes a single Agent performs all cognitive operations sequentially. Parallel literature review, distributed hypothesis generation, and collaborative verification across multiple agents remain future work.

\textbf{(6) Domain concentration.} The skill library and produced paper drafts are strongest in biomedical domains (neurology, ophthalmology, biomechanics), reflecting the developers' domain expertise. Extending to physics, chemistry, or social sciences would require domain-specific knowledge acquisition skills and validation datasets.

\textbf{(7) External evaluation scope.} While the internal benchmark suite demonstrates 100\% pass rate and the golden test suite validates I/O contract compliance, third-party benchmarks (ResearchRubrics, PaperWritingBench full evaluation) remain unaddressed due to the substantial manual effort required for ground-truth construction.

Future work targets three directions: (1) autonomous active inference execution---closing the proposal-to-experiment loop; (2) multi-model epistemic triangulation---running atoms across Model A (deepseek-v4), Model B (claude-sonnet-4), and comparing via VER atom; (3) constitutional evolution---formal amendment procedures that allow the CONSTITUTION.md itself to evolve through structured debate and vote.

\section{Conclusion}

We presented Synthos, a self-evolving cognitive operating system that reifies eight epistemological principles as executable SKILL.md atoms governed by a formal constitutional hierarchy. Over 41 evolution cycles, the system achieved a composite quality score of 0.975, absorbed patterns from 12 external projects, maintained all cognitive atoms at perfect scores, and produced 55+ paper drafts. By demonstrating that philosophical principles can be operationalized as executable cognitive constraints---and that the resulting architecture can deliberately improve itself---Synthos offers a path toward AI research agents that are not only capable but epistemologically rigorous and self-improving.

\begin{thebibliography}{34}

\bibitem{elovic2024gpt}
A. Elovic, ``GPT Researcher,'' 2024. [Online]. Available: \url{https://github.com/assafelovic/gpt-researcher}

\bibitem{sakana2025ai}
Sakana AI, ``The AI Scientist: Towards Fully Automated Open-Ended Scientific Discovery,'' 2025. [Online]. Available: \url{https://github.com/SakanaAI/AI-Scientist}

\bibitem{shao2024storm}
Y. Shao et al., ``Assisting in Writing Wikipedia-like Articles From Scratch with Large Language Models,'' \textit{arXiv:2402.14207}, 2024.

\bibitem{futurehouse2025paperqa2}
Future House, ``PaperQA2: Retrieval-Augmented Generation for Scientific Literature,'' in \textit{Proc. ICLR}, 2025. [Online]. Available: \url{https://github.com/Future-House/paper-qa}

\bibitem{asai2025openscholar}
A. Asai et al., ``OpenScholar: Synthesizing Scientific Literature with Retrieval-Augmented LMs,'' \textit{Nature}, 2025.

\bibitem{hermes2025hermes}
Nous Research, ``Hermes Agent: A Skill-Driven AI Agent Framework,'' 2025. [Online]. Available: \url{https://hermes-agent.nousresearch.com}

\bibitem{khattab2023dspy}
O. Khattab et al., ``DSPy: Compiling Declarative Language Model Calls into Self-Improving Pipelines,'' \textit{arXiv:2310.03714}, 2023.

\bibitem{langchain2024langgraph}
LangChain, ``LangGraph: Building Stateful, Multi-Actor Applications with LLMs,'' 2024. [Online]. Available: \url{https://github.com/langchain-ai/langgraph}

\bibitem{anthropic2025claude}
Anthropic, ``Claude Code: Agentic Coding with Constitutional AI,'' 2025. [Online]. Available: \url{https://github.com/anthropics/claude-code}

\bibitem{skyworkai2026deep}
SkyworkAI, ``Deep Research Agent: A Self-Evolution Protocol and Runtime,'' 2026. [Online]. Available: \url{https://github.com/SkyworkAI/DeepResearchAgent}

\bibitem{orchestra2026airesearch}
Orchestra Research, ``AI Research SKILLs: 90 Skills for Autonomous ML Research,'' 2026. [Online]. Available: \url{https://github.com/Orchestra-Research/AI-research-SKILLs}

\bibitem{yang2026synthos}
X. Yang, ``Synthos: Autonomous Evolving Research OS,'' 2026. [Online]. Available: \url{https://github.com/yakeworld/Synthos}

\bibitem{yang2026synthosconst}
X. Yang, ``Synthos Constitution v5.0 -- Immutable Principles for AI Cognitive Systems,'' 2026.

\bibitem{song2026paperorchestra}
Song et al., ``PaperOrchestra: Orchestrating Scientific Paper Writing with Multi-Agent Collaboration,'' \textit{arXiv:2604.05018}, 2026.

\bibitem{haiyang2025aris}
Haiyang Diiing, ``ARIS: Autonomous Research Iteration System,'' 2025. [Online]. Available: \url{https://github.com/HaiyangDiiing/aris-cli}

\bibitem{gepa2025gepa}
GEPA Team, ``GEPA: General Evolutionary Planning for Agents,'' 2025. [Online]. Available: \url{https://github.com/gepa-ai/gepa}

\bibitem{openai2025agents}
OpenAI, ``OpenAI Agents SDK,'' 2025. [Online]. Available: \url{https://github.com/openai/openai-agents-python}

\bibitem{aiming2025claw}
Aiming Lab, ``AutoResearchClaw: Autonomous Research Agent,'' 2025. [Online]. Available: \url{https://github.com/aiming-lab/AutoResearchClaw}

\bibitem{vo2025kilo}
V. D. Loc, ``KILO-KIT: CRISP-DM Templates for Knowledge-Driven Research,'' 2025. [Online]. Available: \url{https://github.com/VoDaiLocz/KILO-KIT}

\bibitem{ars2025academic}
ARS Team, ``Academic Research Skills v3.7,'' 2026. [Online]. Available: \url{https://github.com/Imbad0202/academic-research-skills}

\bibitem{yuan2025nature}
Yuan1z0825, ``nature-skills: Academic Task Skills for Language Agents,'' 2025. [Online]. Available: \url{https://github.com/Yuan1z0825/nature-skills}

\bibitem{autox2025autor}
AutoX-AI-Labs, ``AutoR: A Human-Centered Research OS,'' 2025. [Online]. Available: \url{https://github.com/AutoX-AI-Labs/AutoR}

\bibitem{friston2010free}
K. Friston, ``The Free-Energy Principle: A Unified Brain Theory?'' \textit{Nature Reviews Neuroscience}, vol.~11, pp.~127--138, 2010.

\bibitem{bengio2013representation}
Y. Bengio, A. Courville, and P. Vincent, ``Representation Learning: A Review and New Perspectives,'' \textit{IEEE Trans. Pattern Analysis and Machine Intelligence}, vol.~35, no.~8, pp.~1798--1828, 2013.

\bibitem{vaswani2017attention}
A. Vaswani et al., ``Attention Is All You Need,'' in \textit{Proc. NeurIPS}, 2017.

\bibitem{brown2020language}
T. B. Brown et al., ``Language Models are Few-Shot Learners,'' in \textit{Proc. NeurIPS}, 2020.

\bibitem{wei2022chain}
J. Wei et al., ``Chain-of-Thought Prompting Elicits Reasoning in Large Language Models,'' in \textit{Proc. NeurIPS}, 2022.

\bibitem{wang2023selfconsistency}
X. Wang et al., ``Self-Consistency Improves Chain of Thought Reasoning in Language Models,'' in \textit{Proc. ICLR}, 2023.

\bibitem{yao2023tree}
S. Yao et al., ``Tree of Thoughts: Deliberate Problem Solving with Large Language Models,'' in \textit{Proc. NeurIPS}, 2023.

\bibitem{lewis2020retrieval}
P. Lewis et al., ``Retrieval-Augmented Generation for Knowledge-Intensive NLP Tasks,'' in \textit{Proc. NeurIPS}, 2020.

\bibitem{chase2022langchain}
H. Chase, ``LangChain: Building Applications with LLMs through Composability,'' 2022. [Online]. Available: \url{https://github.com/langchain-ai/langchain}

\bibitem{sumers2023cognitive}
T. Sumers et al., ``Cognitive Architectures for Language Agents,'' \textit{arXiv:2309.02427}, 2023.

\bibitem{park2023generative}
J. S. Park et al., ``Generative Agents: Interactive Simulacra of Human Behavior,'' in \textit{Proc. UIST}, 2023.

\bibitem{gao2024scientific}
Y. Gao et al., ``Scientific Discovery in the Age of Large Language Models,'' \textit{Nature Reviews Physics}, vol.~6, pp.~325--340, 2024.

\bibitem{wang2024surveyagent}
L. Wang et al., ``A Survey on Large Language Model based Autonomous Agents,'' \textit{arXiv:2308.11432}, 2024.

\bibitem{xi2023rise}
Z. Xi et al., ``The Rise and Potential of Large Language Model Based Agents: A Survey,'' \textit{arXiv:2309.07864}, 2023.

\end{thebibliography}

\end{document}
